\chapter{Sample Spaces, Events, and the Axioms}

So far we have spoken of probability the way most people first meet it: a vague
feeling that some things are more likely than others, sharpened by counting in
the special case where outcomes are ``equally likely.'' That is a fine place to
start, but it is not a foundation. If a friend says ``the probability is
$0.7$,'' what \emph{object} is being assigned the number $0.7$? What rules must
those numbers obey so that they hang together into a consistent picture instead
of contradicting one another? And how do we talk about probability when the
outcomes are \emph{not} equally likely, or when there are infinitely many of
them, or uncountably many?

This chapter builds the stage on which all the rest of the book performs. We
will do three things. First, we cast every random situation as a \emph{set}: a
sample space of possible outcomes, with events as subsets. Second, we write down
the three rules --- Kolmogorov's axioms --- that any honest assignment of
probabilities must satisfy, and we argue carefully \emph{why exactly these three}
and not some other list. Third, we squeeze those three rules for everything they
are worth, deriving the everyday tools (complements, the union bound,
monotonicity, inclusion--exclusion) as \emph{theorems} rather than as separate
articles of faith. Along the way we meet the gentle idea of an event algebra,
which quietly becomes the $\sigma$-algebra of the measure-theoretic chapters at
the end of the book.

The payoff is enormous and a little surprising: from a list of three modest-looking
rules, the entire edifice of probability follows by pure logic. That is the
mathematician's dream, realized for chance.

\section{Outcomes, sample spaces, and events}
\label{ch03:sets}

\begin{intuitionbox}
Before any formalism: imagine you are about to do something whose result you
cannot predict --- flip a coin, roll dice, draw a card, measure tomorrow's
rainfall, watch a stock until noon. Probability begins by listing
\emph{everything that could happen}. That master list is the sample space.
Anything you might later care to bet on --- ``the coin shows heads,'' ``the dice
sum to seven,'' ``it rains more than an inch'' --- is just a \emph{portion} of
that list, a subset. The whole subject is, at bottom, the study of how to put
sensible numbers on subsets.
\end{intuitionbox}

\begin{definition}[Sample space and outcomes]
\label{ch03:def-sample-space}
A \textbf{random experiment} is any process whose outcome is uncertain in
advance. The \textbf{sample space}, written $\Omega$ (capital Greek omega), is
the set of all possible outcomes of the experiment. A single element
$\omega \in \Omega$ is called an \textbf{outcome} or a \textbf{sample point}.
\end{definition}

The only requirement on $\Omega$ is that it be \emph{exhaustive} (every possible
result is listed) and that its elements be \emph{mutually exclusive} (exactly one
$\omega$ actually occurs on any given run). These two conditions are the entire
content of ``list everything that could happen, with no overlaps and no gaps.''

\begin{example}[Three classic sample spaces]
\label{ch03:ex-three-spaces}
\begin{itemize}
\item One coin flip: $\Omega = \set{H, T}$.
\item One die roll: $\Omega = \set{1,2,3,4,5,6}$.
\item Two distinguishable dice: $\Omega = \set{(i,j) : i,j \in \set{1,\dots,6}}$,
a set of $36$ ordered pairs.
\end{itemize}
\end{example}

A handy way to \emph{generate} a sample space for a multi-stage experiment is to
grow it stage by stage as a tree: each level branches over the possibilities at
that step, and the leaves are the outcomes. For three coin flips the tree has
$2 \times 2 \times 2 = 8$ leaves, which is exactly $\abs{\Omega}$ for
$\Omega = \set{H,T}^3$ (Figure~\ref{ch03:fig-tree}).

\begin{figure}[H]
\centering
\begin{tikzpicture}[
  level distance=14mm,
  level 1/.style={sibling distance=42mm},
  level 2/.style={sibling distance=21mm},
  level 3/.style={sibling distance=10.5mm},
  edge from parent/.style={draw, gray!70, -{Latex[length=1.6mm]}},
  every node/.style={font=\footnotesize},
]
\node[circle, draw, fill=blue!5, inner sep=1.5pt] {start}
  child {node[circle, draw, fill=sourceblue!12, inner sep=1pt] {H}
    child {node[circle, draw, inner sep=1pt] {H}
      child {node {HHH}}
      child {node {HHT}}}
    child {node[circle, draw, inner sep=1pt] {T}
      child {node {HTH}}
      child {node {HTT}}}}
  child {node[circle, draw, fill=targetgreen!12, inner sep=1pt] {T}
    child {node[circle, draw, inner sep=1pt] {H}
      child {node {THH}}
      child {node {THT}}}
    child {node[circle, draw, inner sep=1pt] {T}
      child {node {TTH}}
      child {node {TTT}}}};
\end{tikzpicture}
\caption{A tree for three flips of a coin. Each level branches over the result of
one flip; the eight leaves are exactly the outcomes of
$\Omega=\set{H,T}^3$. Trees make the multiplication of possibilities visible and
are the natural scaffold for staged experiments.}
\label{ch03:fig-tree}
\end{figure}

\begin{definition}[Event]
\label{ch03:def-event}
An \textbf{event} is a subset $A \subseteq \Omega$. We say the event $A$
\textbf{occurs} on a particular run of the experiment if the realized outcome
$\omega$ lies in $A$. The whole space $\Omega$ is the \textbf{sure event} (it
always occurs); the empty set $\varnothing$ is the \textbf{impossible event} (it
never occurs). A one-element event $\set{\omega}$ is sometimes called an
\textbf{elementary event}.
\end{definition}

This identification --- \emph{event = subset} --- is the quiet hinge on which
the whole theory turns, because it lets us translate ordinary spoken language
about chance into the precise grammar of set theory. The translation table is
worth memorizing not as a chore but as a habit of seeing.

\begin{focusbox}
\textbf{The dictionary between English and sets.} Let $A$ and $B$ be events.
\begin{center}
\begin{tabular}{@{}ll@{}}
\toprule
\textbf{English} & \textbf{Set notation} \\
\midrule
``$A$ does not occur'' & $A^c = \Omega \setminus A$ (complement) \\
``$A$ or $B$ occurs'' (inclusive) & $A \cup B$ (union) \\
``both $A$ and $B$ occur'' & $A \cap B$ (intersection) \\
``$A$ but not $B$'' & $A \cap B^c = A \setminus B$ \\
``$A$ implies $B$'' & $A \subseteq B$ \\
``$A$ and $B$ cannot both occur'' & $A \cap B = \varnothing$ (disjoint) \\
``at least one of $A_1, A_2, \dots$'' & $\bigcup_i A_i$ \\
``all of $A_1, A_2, \dots$'' & $\bigcap_i A_i$ \\
\bottomrule
\end{tabular}
\end{center}
\end{focusbox}

Two events with empty intersection, $A \cap B = \varnothing$, are called
\textbf{disjoint} or \textbf{mutually exclusive}: they cannot occur together. A
collection of events that are pairwise disjoint and whose union is all of
$\Omega$ is a \textbf{partition} --- it carves the sample space into
non-overlapping pieces that together cover everything. Partitions will become one
of our sharpest tools when we reach the law of total probability.

\begin{example}[Reading events off two dice]
\label{ch03:ex-two-dice-events}
On $\Omega = \set{(i,j)}$ for two dice, consider:
\begin{itemize}
\item $A = $ ``the sum is $7$'' $= \set{(1,6),(2,5),(3,4),(4,3),(5,2),(6,1)}$,
which has six outcomes.
\item $B = $ ``the first die is even'' $= \set{(i,j): i \in \set{2,4,6}}$, which
has $18$ outcomes.
\item $A \cap B = \set{(2,5),(4,3),(6,1)}$: sum seven \emph{and} first die even.
\item $A^c = $ ``the sum is not $7$,'' the remaining $30$ outcomes.
\end{itemize}
Notice we have not yet assigned any probabilities. We have only described
\emph{which outcomes belong to which events}. Geometry first, numbers later.
\end{example}

\begin{historybox}
The word \emph{event} (\'ev\'enement) and the deliberate habit of treating chance
problems as questions about \emph{sets of cases} go back to the founding
correspondence between Blaise Pascal and Pierre de Fermat in 1654, prompted by
gambling questions posed by the Chevalier de M\'er\'e. But it was the
nineteenth-century rise of set theory under Georg Cantor, and then the
twentieth-century work of \'Emile Borel and Henri Lebesgue on \emph{measure} ---
on assigning sizes to subsets of the line --- that supplied the language in which
``event = subset'' could be made completely rigorous. Andrey Kolmogorov's 1933
monograph \emph{Grundbegriffe der Wahrscheinlichkeitsrechnung} (``Foundations of
the Theory of Probability'') brought these threads together and is the direct
ancestor of every definition in this chapter.
\end{historybox}

\subsection{Sample spaces come in three sizes}
\label{ch03:three-sizes}

The character of a probability problem is set largely by how big $\Omega$ is, and
there are exactly three regimes worth distinguishing.

\begin{itemize}
\item \textbf{Finite.} $\Omega$ has finitely many outcomes: coins, dice, cards, a
poll of $1000$ voters. Here we can in principle list every outcome and the only
hard part is counting (Chapter on combinatorics).
\item \textbf{Countably infinite.} $\Omega$ is infinite but can be listed
$\omega_1, \omega_2, \omega_3, \dots$, like the natural numbers. Example: flip a
coin until the first head and record how many flips it took, giving
$\Omega = \set{1, 2, 3, \dots}$. We can still assign a probability to each
outcome individually and add them up --- but now the sum is an infinite series,
so convergence enters the story.
\item \textbf{Uncountable.} $\Omega$ is a continuum, like an interval of the real
line: tomorrow's temperature, the exact time a bus arrives, a dart's landing
point. Now the outcomes cannot be listed, and --- this is the subtle part ---
typically \emph{every individual outcome has probability zero}, yet whole
intervals have positive probability. Here probability becomes \emph{area under a
curve}, and calculus takes over.
\end{itemize}

\begin{connectionbox}
The jump from countable to uncountable is precisely the jump from \emph{sums} to
\emph{integrals}. In a discrete problem the probability of an event is a sum of
the probabilities of its outcomes, $\P(A) = \sum_{\omega \in A} \P(\set{\omega})$.
In a continuous problem it becomes $\P(A) = \int_A f(x)\,\dd x$ for a density
$f$. The whole machinery of the calculus --- limits, areas, the fundamental
theorem --- is exactly what we will borrow to handle uncountable sample spaces in
the continuous chapters. If the symbol $\int_A f\,\dd x$ feels rusty, the
companion calculus book is the place to refresh; here we only need to believe
that ``total area $=1$'' is the continuous echo of ``probabilities sum to one.''
\end{connectionbox}

\begin{trapbox}
\textbf{The sample space is a modeling choice, not a fact of nature.} For two
dice you could use the $36$ ordered pairs, or you could use the $11$ possible
sums $\set{2,3,\dots,12}$, or the $21$ unordered pairs. All three are legitimate
sample spaces for ``roll two dice''! They differ in what they let you talk about
and in whether the outcomes are equally likely. The $36$ ordered pairs are
equally likely; the $11$ sums are emphatically \emph{not} (a seven is six times as
likely as a two). Choosing $\Omega$ well --- usually so that the outcomes are
equally likely --- is half the art of solving a probability problem.
\end{trapbox}

\section{The Kolmogorov axioms}
\label{ch03:axioms}

We now ask the central question of the chapter: what is a probability, really?
We want a function $\P$ that takes an event and returns a number measuring how
likely it is. Rather than try to define ``likely'' from scratch --- a swamp of
philosophy --- we follow Kolmogorov's brilliant strategy: we simply write down
the handful of properties any reasonable likelihood-measure \emph{must} have, and
then declare \emph{anything} satisfying them to be a probability. This is the
axiomatic method, the same move Euclid made for geometry.

\begin{intuitionbox}
What do we already believe about likelihood, before any theory? Three things,
really. \textbf{(1)} A likelihood is never negative --- it makes no sense for an
event to be ``minus thirty percent'' likely. \textbf{(2)} Something is certain to
happen, so the whole sample space should carry the maximal likelihood, which we
normalize to $1$. \textbf{(3)} If two events cannot both occur, the chance that
\emph{one or the other} happens is just the sum of their chances --- likelihoods
of separate, non-overlapping possibilities add up. That is the entire intuition.
The axioms below are just these three sentences written carefully.
\end{intuitionbox}

\begin{definition}[Probability measure; the Kolmogorov axioms]
\label{ch03:def-axioms}
Let $\Omega$ be a sample space and let $\mathcal{F}$ be a collection of events
(subsets of $\Omega$) on which we wish to define probabilities. A
\textbf{probability measure} is a function $\P : \mathcal{F} \to \R$ satisfying:
\begin{description}
\item[Axiom 1 (Nonnegativity).] $\P(A) \ge 0$ for every event $A$.
\item[Axiom 2 (Normalization).] $\P(\Omega) = 1$.
\item[Axiom 3 (Countable additivity).] For any \emph{countable} collection of
\emph{pairwise disjoint} events $A_1, A_2, A_3, \dots$ (so $A_i \cap A_j =
\varnothing$ whenever $i \ne j$),
\[
\P\!\left(\bigcup_{i=1}^{\infty} A_i\right) = \sum_{i=1}^{\infty} \P(A_i).
\]
\end{description}
The triple $(\Omega, \mathcal{F}, \P)$ is called a \textbf{probability space}.
\end{definition}

That is the whole foundation. Everything else in probability theory --- the
normal distribution, the central limit theorem, Markov chains, Brownian motion ---
is a logical consequence of these three lines together with definitions. It is
worth pausing on how little we assumed and how much will follow.

\begin{whybox}
\textbf{Why exactly these three, and why countable additivity in particular?}

\emph{Axiom 1} is forced: probabilities are sizes, and sizes are not negative.

\emph{Axiom 2} is a normalization, a choice of units. We could have used a scale
from $0$ to $100$ (percent) or $0$ to $1$. Setting $\P(\Omega) = 1$ makes the
arithmetic clean: probabilities become fractions of the total certainty, and a
probability is literally ``the share of the sure thing that this event claims.''

\emph{Axiom 3} is the real engine, and the choice of \emph{countable} (rather
than merely \emph{finite}) additivity is the one genuinely deep decision. Finite
additivity --- $\P(A \cup B) = \P(A) + \P(B)$ for disjoint $A, B$ --- is obvious
and would be agreed to by anyone. The strengthening to countably many sets is
what lets us take \emph{limits}, and limits are the lifeblood of analysis. Want
to compute the probability that a coin eventually shows heads? That is a union of
infinitely many disjoint events (``first head on flip $1$,'' ``first head on flip
$2$,'' \dots), and only countable additivity lets us add their probabilities to
get the answer. Without it, the continuity properties at the end of this chapter
fail, and with them every limit theorem in the book. Kolmogorov chose countable
additivity precisely because it is the weakest assumption strong enough to make
probability a branch of \emph{measure theory}, where the full power of calculus
becomes available.
\end{whybox}

\begin{happeningbox}
\textbf{Probability is normalized area.} Picture $\Omega$ as a region in the
plane with total area $1$. An event $A$ is a sub-region, and $\P(A)$ is simply
\emph{its area}. Axiom 1 says regions do not have negative area. Axiom 2 says the
whole board has area $1$. Axiom 3 says that if you chop a region into
non-overlapping pieces, the areas of the pieces add up to the area of the whole.
Stated this way the axioms are not mysterious at all --- they are the elementary
school facts about area, transplanted onto the sample space. Everything we prove
below is something you already believe about areas; we are just being careful
enough to be sure. This picture is literally correct: a probability measure
\emph{is} a measure in the sense of Lebesgue, with the single extra condition
that the total measure is $1$.
\end{happeningbox}

\begin{figure}[H]
\centering
\begin{tikzpicture}
  % The sample space as a rectangle of area 1
  \draw[thick, fill=blue!4] (0,0) rectangle (8,4.6);
  \node[anchor=north east] at (8,4.6) {\large $\Omega$};
  % Events A and B as overlapping ellipses
  \begin{scope}
    \clip (0,0) rectangle (8,4.6);
    \fill[sourceblue!35] (2.3,2.3) ellipse (1.9 and 1.4);
    \fill[targetgreen!35] (5.1,2.0) ellipse (1.7 and 1.5);
    \begin{scope}
      \clip (2.3,2.3) ellipse (1.9 and 1.4);
      \fill[structurepurple!55] (5.1,2.0) ellipse (1.7 and 1.5);
    \end{scope}
  \end{scope}
  \draw[thick, sourceblue] (2.3,2.3) ellipse (1.9 and 1.4);
  \draw[thick, targetgreen] (5.1,2.0) ellipse (1.7 and 1.5);
  \node[sourceblue!60!black] at (1.5,3.2) {\large $A$};
  \node[targetgreen!60!black] at (6.2,2.9) {\large $B$};
  \node[structurepurple!60!black, font=\small] at (3.85,1.4) {$A\cap B$};
\end{tikzpicture}
\caption{The area picture. The sample space $\Omega$ is a board of total area
$1$; events are sub-regions and $\P$ measures their area. When $A$ and $B$
overlap, the overlap $A\cap B$ is counted in both $\P(A)$ and $\P(B)$, which is
the seed of inclusion--exclusion.}
\label{ch03:fig-area}
\end{figure}

\begin{formalbox}
A subtle point hides in Definition~\ref{ch03:def-axioms}: we said $\P$ is defined
on a \emph{collection} $\mathcal{F}$ of events, not necessarily on \emph{all}
subsets of $\Omega$. For finite or countable $\Omega$ we may safely take
$\mathcal{F}$ to be the set of all subsets, and you may ignore the distinction.
But for uncountable $\Omega$ --- the real line, say --- it turns out to be
mathematically \emph{impossible} to assign a sensible probability to every single
subset at once. We are forced to restrict $\mathcal{F}$ to a well-behaved family,
a \emph{$\sigma$-algebra}. We preview this in Section~\ref{ch03:algebra} and treat
it in full in the measure-theory chapter. For now, whenever we write $\P(A)$,
trust that $A$ is one of the allowed events.
\end{formalbox}

\section{Basic consequences of the axioms}
\label{ch03:consequences}

Now we harvest. Each result below is a theorem --- it \emph{follows} from the
three axioms by logic alone. We prove them honestly and in an order where each
leans on the previous ones. Treat this section as a worked demonstration that a
tiny axiom set is astonishingly fertile.

\subsection{From countable to finite additivity}
\label{ch03:finite-add}

First a housekeeping step: Axiom 3 is stated for infinitely many sets, but we
usually have only two or three. The finite version is a corollary.

\begin{proposition}[Finite additivity]
\label{ch03:prop-finite-add}
If $A_1, \dots, A_n$ are pairwise disjoint events, then
$\P\!\left(\bigcup_{i=1}^{n} A_i\right) = \sum_{i=1}^{n} \P(A_i)$. In particular,
for disjoint $A, B$: \ $\P(A \cup B) = \P(A) + \P(B)$.
\end{proposition}

\begin{proof}
We first show $\P(\varnothing) = 0$. The empty set is disjoint from itself, and
$\varnothing = \varnothing \cup \varnothing \cup \cdots$. Apply Axiom 3 to the
countable disjoint family $A_1 = A_2 = \cdots = \varnothing$:
$\P(\varnothing) = \sum_{i=1}^\infty \P(\varnothing)$. The only nonnegative
number equal to an infinite sum of copies of itself is $0$, so
$\P(\varnothing) = 0$.

Now take the finite family $A_1, \dots, A_n$ and pad it with empties:
$A_{n+1} = A_{n+2} = \cdots = \varnothing$. This is a countable disjoint family
with the same union, so Axiom 3 gives
\[
\P\!\left(\bigcup_{i=1}^{n} A_i\right)
= \sum_{i=1}^{\infty} \P(A_i)
= \sum_{i=1}^{n} \P(A_i) + \sum_{i>n} \P(\varnothing)
= \sum_{i=1}^{n} \P(A_i),
\]
since the tail terms are all $0$.
\end{proof}

Notice that countable additivity quietly did two jobs: it gave us
$\P(\varnothing)=0$ \emph{and} the finite rule. This is the payoff of choosing the
stronger axiom.

\subsection{The complement rule and the unit interval}
\label{ch03:complement}

\begin{proposition}[Complement rule]
\label{ch03:prop-complement}
For any event $A$, \ $\P(A^c) = 1 - \P(A)$. Consequently $0 \le \P(A) \le 1$.
\end{proposition}

\begin{proof}
The events $A$ and $A^c$ are disjoint and their union is $\Omega$. By finite
additivity and Axiom 2,
\[
1 = \P(\Omega) = \P(A \cup A^c) = \P(A) + \P(A^c),
\]
so $\P(A^c) = 1 - \P(A)$. Rearranging and using Axiom 1 ($\P(A^c) \ge 0$) gives
$\P(A) = 1 - \P(A^c) \le 1$; combined with $\P(A) \ge 0$ we get
$0 \le \P(A) \le 1$.
\end{proof}

So the familiar fact that probabilities live in $[0,1]$ was \emph{not} an axiom
--- only $\P(A)\ge 0$ and $\P(\Omega)=1$ were assumed. The upper bound $\P(A)\le1$
is a theorem we just earned. This is the kind of economy the axiomatic method
buys: assume less, derive more.

\begin{aestheticbox}
The complement rule is the single most useful trick in elementary probability,
because ``not'' is often vastly easier to count than the thing itself. ``What is
the probability that at least one of these ten components fails?'' is a tangle of
overlapping cases. ``What is the probability that \emph{none} fails?'' is a single
clean product. Whenever a problem says ``at least one,'' your reflex should be:
take the complement. We use exactly this move to crack the birthday problem
below.
\end{aestheticbox}

\subsection{Monotonicity and the difference rule}
\label{ch03:monotone}

\begin{proposition}[Monotonicity]
\label{ch03:prop-monotone}
If $A \subseteq B$ then $\P(A) \le \P(B)$, and moreover
$\P(B \setminus A) = \P(B) - \P(A)$.
\end{proposition}

\begin{proof}
Split $B$ into the part inside $A$ and the part outside: $B = A \cup
(B \setminus A)$, where $A$ and $B \setminus A$ are disjoint (the second is
$B \cap A^c$). Finite additivity gives
$\P(B) = \P(A) + \P(B \setminus A)$. Since $\P(B \setminus A) \ge 0$ by Axiom 1,
we get $\P(B) \ge \P(A)$; and rearranging gives the difference formula.
\end{proof}

Monotonicity is the formal statement of an obvious truth: \emph{a bigger event is
at least as likely as a smaller one it contains}. If $A$ implies $B$ (every way
$A$ can happen is also a way $B$ happens), then $B$ cannot be less likely than
$A$. Many ``paradoxes'' in probability are really violations of monotonicity that
a careless verbal description sneaks past you --- the conjunction fallacy, where
people rate ``Linda is a feminist bank teller'' as more probable than ``Linda is
a bank teller,'' is exactly a failure to notice $A \cap B \subseteq B$.

\subsection{The inclusion--exclusion principle}
\label{ch03:incl-excl}

When two events overlap, we cannot simply add their probabilities --- we would
count the overlap twice (look again at Figure~\ref{ch03:fig-area}). The fix is to
subtract the double-counted part.

\begin{theorem}[Inclusion--exclusion, two events]
\label{ch03:thm-ie2}
For any events $A, B$,
\[
\P(A \cup B) = \P(A) + \P(B) - \P(A \cap B).
\]
\end{theorem}

\begin{proof}
We decompose $A \cup B$ into three disjoint pieces: the part in $A$ only, the
part in $B$ only, and the part in both:
\[
A \cup B = (A \setminus B) \;\cup\; (B \setminus A) \;\cup\; (A \cap B),
\]
a disjoint union. By finite additivity,
\[
\P(A \cup B) = \P(A \setminus B) + \P(B \setminus A) + \P(A \cap B).
\]
Now use the difference rule (Proposition~\ref{ch03:prop-monotone}) on the pieces of
$A$ and of $B$: writing $A = (A\setminus B) \cup (A \cap B)$ disjointly gives
$\P(A \setminus B) = \P(A) - \P(A \cap B)$, and likewise
$\P(B \setminus A) = \P(B) - \P(A \cap B)$. Substituting,
\[
\P(A \cup B) = \big[\P(A) - \P(A\cap B)\big] + \big[\P(B) - \P(A\cap B)\big] +
\P(A \cap B) = \P(A) + \P(B) - \P(A \cap B). \qedhere
\]
\end{proof}

\begin{theorem}[Inclusion--exclusion, three events]
\label{ch03:thm-ie3}
For any events $A, B, C$,
\[
\P(A \cup B \cup C) = \P(A) + \P(B) + \P(C) - \P(A\cap B) - \P(A \cap C) -
\P(B \cap C) + \P(A \cap B \cap C).
\]
\end{theorem}

The pattern continues for $n$ events: add all single probabilities, subtract all
pairwise overlaps, add back all triple overlaps, subtract the quadruples, and so
on, alternating signs. The general statement is
\[
\P\!\left(\bigcup_{i=1}^n A_i\right) = \sum_{k=1}^{n} (-1)^{k+1}
\!\!\sum_{1 \le i_1 < \cdots < i_k \le n}\!\! \P(A_{i_1} \cap \cdots \cap A_{i_k}).
\]
We will derive the general version cleanly in a later chapter using indicator
random variables, where it falls out in one line; for now the two- and
three-event cases are all we need.

\begin{example}[Counting with inclusion--exclusion]
\label{ch03:ex-ie}
In a class of $100$ students, $40$ take French, $30$ take German, and $12$ take
both. A student is chosen uniformly at random. What is the probability the
student takes at least one of the two languages?

Let $F$ and $G$ be the events of taking French and German. Then $\P(F) = 0.40$,
$\P(G) = 0.30$, and $\P(F \cap G) = 0.12$. By Theorem~\ref{ch03:thm-ie2},
\[
\P(F \cup G) = 0.40 + 0.30 - 0.12 = 0.58.
\]
Had we forgotten to subtract the overlap, we would have wrongly reported
$0.70$, double-counting the $12$ students who take both. The probability of
taking \emph{neither} language is then $1 - 0.58 = 0.42$ by the complement rule.
\end{example}

\subsection{The union bound (Boole's inequality)}
\label{ch03:union-bound}

Inclusion--exclusion is exact but, for many events, a nightmare of alternating
terms. Very often we do not need the exact value of $\P(\bigcup A_i)$ --- we just
need to know it is \emph{small}. Then one of the most beloved tools in all of
probability comes to the rescue.

\begin{theorem}[Union bound / Boole's inequality]
\label{ch03:thm-union-bound}
For any events $A_1, A_2, A_3, \dots$ (finitely or countably many, disjoint or
not),
\[
\P\!\left(\bigcup_{i} A_i\right) \le \sum_{i} \P(A_i).
\]
\end{theorem}

\begin{proof}
The trick is \emph{disjointification}: replace the $A_i$ by disjoint pieces with
the same union. Define
\[
B_1 = A_1, \qquad B_i = A_i \setminus (A_1 \cup \cdots \cup A_{i-1}) \ \text{for } i \ge 2.
\]
Each $B_i$ is the ``new'' part of $A_i$ not already covered by earlier sets. By
construction the $B_i$ are pairwise disjoint, $B_i \subseteq A_i$, and
$\bigcup_i B_i = \bigcup_i A_i$. Then, using countable additivity on the disjoint
$B_i$ and monotonicity ($B_i \subseteq A_i$ so $\P(B_i) \le \P(A_i)$),
\[
\P\!\left(\bigcup_i A_i\right) = \P\!\left(\bigcup_i B_i\right) = \sum_i \P(B_i)
\le \sum_i \P(A_i). \qedhere
\]
\end{proof}

\begin{intuitionbox}
Why is the union bound true ``morally''? Because when you add up
$\sum_i \P(A_i)$, any outcome lying in several of the $A_i$ gets counted multiple
times, whereas in $\P(\bigcup A_i)$ it is counted once. Overcounting can only
inflate the sum, never shrink it --- so the sum is an \emph{upper bound}. The
inequality is tight (an equality) exactly when the events are disjoint, in which
case there is no overcounting and the union bound reduces to additivity.
\end{intuitionbox}

\begin{applicationbox}
\textbf{The union bound is the workhorse of computer science and statistics.}
Suppose you run $20$ different hypothesis tests, each with a $5\%$ chance of a
false alarm under the null. What is the chance of \emph{at least one} false alarm
across all $20$? Exact computation requires knowing the dependence structure, but
the union bound gives an instant, assumption-free ceiling:
$\P(\text{any false alarm}) \le 20 \times 0.05 = 1.0$ --- here useless, but for
$5$ tests it gives $\le 0.25$, a real guarantee. This is the logic behind the
\emph{Bonferroni correction}: to keep the overall error rate below $\alpha$ across
$m$ tests, test each at level $\alpha/m$, since
$m \cdot (\alpha/m) = \alpha$. The same one-line bound underlies generalization
guarantees in machine learning (bounding the chance that \emph{some} hypothesis
in a class looks good by luck) and tail bounds for the maximum of many random
quantities. A great deal of theoretical computer science is, at heart, the
sustained and clever application of this single inequality.
\end{applicationbox}

\section{Equally likely outcomes}
\label{ch03:equally-likely}

We now connect the abstract axioms back to the concrete counting you already
know. In the simplest and oldest model of chance, every outcome is just as likely
as every other. The axioms then force a beautiful, rigid conclusion.

\begin{proposition}[The classical / equally-likely model]
\label{ch03:prop-classical}
Suppose $\Omega$ is finite with $\abs{\Omega} = N$ outcomes, and that all
outcomes are equally likely. Then each outcome has probability $1/N$, and for any
event $A$,
\[
\P(A) = \frac{\abs{A}}{\abs{\Omega}} = \frac{\text{number of outcomes in } A}{\text{total number of outcomes}}.
\]
\end{proposition}

\begin{proof}
Let $p$ be the common probability of each outcome. The $N$ elementary events
$\set{\omega}$ are disjoint and their union is $\Omega$, so by finite additivity
and normalization, $1 = \P(\Omega) = \sum_{\omega} \P(\set{\omega}) = Np$. Hence
$p = 1/N$. For a general event $A$ with $\abs{A} = k$ outcomes, additivity again
gives $\P(A) = \sum_{\omega \in A} \P(\set{\omega}) = k \cdot \frac{1}{N} =
\abs{A}/N$.
\end{proof}

This is the formula --- ``favorable over total'' --- that probability is often
introduced with, and now you see it is not a definition but a \emph{theorem},
valid \emph{only} when outcomes are equally likely. In that regime, computing
probabilities reduces entirely to \emph{counting}, which is why combinatorics is
the constant companion of elementary probability.

\begin{historybox}
This ``favorable over total'' rule is the original, pre-axiomatic definition of
probability, codified by Pierre-Simon Laplace in his \emph{Th\'eorie analytique
des probabilit\'es} (1812) as the ratio of favorable to ``equally possible''
cases. For a century it \emph{was} the definition. Its fatal weakness is the
circularity buried in ``equally possible'' --- you need a notion of equal
likelihood to define likelihood --- and its uselessness when outcomes are not
symmetric or when there are infinitely many. Kolmogorov's axioms (1933) dethroned
it: now ``equally likely'' is just one special \emph{assumption} you may or may
not make, and the rule of Proposition~\ref{ch03:prop-classical} is what that
assumption implies. The classical model became a corollary instead of a creed.
\end{historybox}

\begin{example}[Two dice, carefully]
\label{ch03:ex-dice-sum}
Roll two fair dice. Use the equally-likely space $\Omega$ of $36$ ordered pairs.
What is $\P(\text{sum} = 7)$? From Example~\ref{ch03:ex-two-dice-events}, the event
has six outcomes, so $\P = 6/36 = 1/6$. What about $\P(\text{sum} = 2)$? Only
$(1,1)$ works, so $\P = 1/36$. This is exactly why the alternative ``sums''
sample space $\set{2,\dots,12}$ is \emph{not} equally likely, and why choosing the
$36$-pair space was the right modeling move (recall the trap box of
Section~\ref{ch03:three-sizes}).
\end{example}

\begin{figure}[H]
\centering
\begin{tikzpicture}
\begin{axis}[
  width=12cm, height=6cm,
  ybar, bar width=14pt,
  xlabel={sum of two fair dice}, ylabel={probability},
  xtick={2,3,4,5,6,7,8,9,10,11,12},
  ymin=0, ymax=0.19,
  ytick={0,0.05,0.1,0.15},
  enlarge x limits=0.06,
  axis lines=left,
  every axis plot/.append style={fill=distblue!65, draw=distblue!90!black},
]
\addplot coordinates {
  (2,0.0278) (3,0.0556) (4,0.0833)
  (5,0.1111) (6,0.1389) (7,0.1667)
  (8,0.1389) (9,0.1111) (10,0.0833)
  (11,0.0556) (12,0.0278)
};
\end{axis}
\end{tikzpicture}
\caption{The probability mass over the eleven possible sums of two fair dice.
Each of the $36$ ordered pairs is equally likely, but the \emph{sums} are not:
seven (six ways, probability $6/36$) is the peak, twos and twelves (one way each)
the valleys. The triangular shape is the first hint of how adding independent
things builds up toward a bell.}
\label{ch03:fig-dice}
\end{figure}

\begin{paradoxbox}
\textbf{The birthday problem.} In a room of just $23$ people, the probability
that two share a birthday is already better than even --- about $50.7\%$. With
$50$ people it rockets to $97\%$. People find this shocking because they
unconsciously ask the wrong question (``does someone share \emph{my} birthday?'')
instead of the right one (``does \emph{any} pair share a birthday?''). The number
of \emph{pairs} among $23$ people is $\binom{23}{2} = 253$, and it is those $253$
chances, not $23$, that drive the collision.

Here is the clean computation, using the complement trick of
Section~\ref{ch03:complement}. Model the $n$ birthdays as a uniform outcome in
$\Omega = \set{1,\dots,365}^n$, all $365^n$ sequences equally likely (ignoring
leap years and seasonal effects). Let $A$ be ``at least one shared birthday.''
Directly, $A$ is a mess of overlapping coincidences. Its complement $A^c$ =
``all $n$ birthdays distinct'' is a clean count: the first person is free
($365$ choices), the second must dodge one date ($364$), the third must dodge two
($363$), and so on. By the equally-likely formula,
\[
\P(A^c) = \frac{365 \cdot 364 \cdots (365-n+1)}{365^{\,n}}, \qquad
\P(A) = 1 - \P(A^c).
\]
For $n=23$ this gives $\P(A^c) \approx 0.4927$, so $\P(A) \approx 0.5073$. The
lesson generalizes far beyond birthdays: coincidences in \emph{any} large set of
\emph{pairwise} comparisons --- hash collisions in computing, matching DNA
profiles, ``amazing'' shared traits at a party --- are far more common than raw
intuition suggests, because the number of pairs grows like $n^2/2$.
\end{paradoxbox}

\begin{figure}[H]
\centering
\begin{tikzpicture}
\begin{axis}[
  width=12cm, height=6cm,
  xlabel={number of people $n$}, ylabel={$\P(\text{a shared birthday})$},
  xmin=0, xmax=60, ymin=0, ymax=1.02,
  xtick={0,10,20,23,30,40,50,60},
  ytick={0,0.25,0.5,0.75,1.0},
  axis lines=left, grid=major, grid style={gray!18},
  samples=120, domain=1:60,
]
% P(match) approximated smoothly by 1 - exp(-n(n-1)/730)
\addplot[densityred, very thick, smooth] {1 - exp(-x*(x-1)/730)};
\addplot[dashed, gray] coordinates {(0,0.5) (23,0.5)};
\addplot[dashed, gray] coordinates {(23,0) (23,0.5)};
\node[anchor=south, font=\footnotesize] at (axis cs:23,0.50) {$n=23$};
\end{axis}
\end{tikzpicture}
\caption{The birthday curve. The probability of a shared birthday crosses $1/2$
at just $n=23$ people (dashed lines) and is near-certain by $n=50$. The S-shape
reflects the quadratic growth in the number of pairs. (Drawn via the accurate
approximation $\P(A)\approx 1-e^{-n(n-1)/730}$.)}
\label{ch03:fig-birthday}
\end{figure}

\section{Continuity of probability}
\label{ch03:continuity}

We promised that countable additivity buys us the ability to take limits. Here is
the cash payment. The idea: if a sequence of events \emph{grows} steadily toward
a limiting event, their probabilities should grow toward the limit's probability
--- and likewise for shrinking sequences. This is so natural you might assume it
needs no proof, but in fact it is \emph{equivalent} to countable additivity, and
it is the property a merely finitely-additive theory lacks.

\begin{definition}[Increasing and decreasing sequences of events]
\label{ch03:def-mono-seq}
A sequence $A_1, A_2, \dots$ is \textbf{increasing} if $A_1 \subseteq A_2
\subseteq A_3 \subseteq \cdots$; its limit is $A = \bigcup_{n} A_n$. It is
\textbf{decreasing} if $A_1 \supseteq A_2 \supseteq \cdots$; its limit is
$A = \bigcap_{n} A_n$.
\end{definition}

\begin{theorem}[Continuity of probability]
\label{ch03:thm-continuity}
If $A_n \uparrow A$ (increasing, with union $A$), then $\P(A_n) \to \P(A)$. If
$A_n \downarrow A$ (decreasing, with intersection $A$), then likewise
$\P(A_n) \to \P(A)$.
\end{theorem}

\begin{proof}
\emph{Increasing case.} Turn the nested sequence into disjoint shells. Set
$B_1 = A_1$ and $B_n = A_n \setminus A_{n-1}$ for $n \ge 2$. The $B_n$ are
pairwise disjoint, $\bigcup_{k=1}^n B_k = A_n$, and $\bigcup_{k=1}^\infty B_k =
A$. By countable additivity and then the definition of an infinite sum as the
limit of partial sums,
\[
\P(A) = \sum_{k=1}^{\infty} \P(B_k) = \lim_{n\to\infty} \sum_{k=1}^{n} \P(B_k)
= \lim_{n\to\infty} \P(A_n),
\]
where the last equality uses finite additivity, $\sum_{k=1}^n \P(B_k) =
\P\!\big(\bigcup_{k=1}^n B_k\big) = \P(A_n)$.

\emph{Decreasing case.} Take complements to reduce to the increasing case. If
$A_n \downarrow A$ then $A_n^c \uparrow A^c$, so the increasing case gives
$\P(A_n^c) \to \P(A^c)$. By the complement rule $\P(A_n^c) = 1 - \P(A_n)$ and
$\P(A^c) = 1 - \P(A)$, hence $1 - \P(A_n) \to 1 - \P(A)$, i.e.\
$\P(A_n) \to \P(A)$.
\end{proof}

\begin{happeningbox}
\textbf{What continuity really lets you do is reason about the infinite future.}
``The coin eventually shows heads,'' ``the random walk returns to the origin,''
``the process never escapes this region'' --- all of these are statements about
\emph{limits} of events, not about any single finite stage. Continuity is the
bridge that lets us compute the probability of such limiting events as the limit
of probabilities we \emph{can} compute at each finite stage. Strip out countable
additivity, and that bridge collapses: you could no longer say that the chance of
``eventually heads'' is the limit of ``heads within the first $n$ flips.'' Every
asymptotic theorem in this book --- the laws of large numbers, recurrence of
random walks, the Borel--Cantelli lemmas --- crosses this bridge.
\end{happeningbox}

\begin{example}[Eventually heads]
\label{ch03:ex-eventually}
Flip a fair coin repeatedly. Let $A_n = $ ``at least one head in the first $n$
flips.'' These events increase, $A_n \uparrow A$ where $A = $ ``a head appears at
some point.'' We can compute each finite term with the complement rule:
$\P(A_n) = 1 - \P(\text{all tails in } n) = 1 - (1/2)^n$. By continuity,
\[
\P(A) = \lim_{n \to \infty} \P(A_n) = \lim_{n\to\infty}\big(1 - 2^{-n}\big) = 1.
\]
A head appears eventually with probability $1$ --- certain, in the
probabilistic sense, even though no individual flip is certain and the
all-tails-forever outcome is not logically \emph{impossible}. It is merely
\emph{null}: it has probability zero. This distinction between ``impossible'' and
``probability zero'' is one of the most important refinements the axioms give
us, and it only makes sense because of continuity.
\end{example}

\begin{figure}[H]
\centering
\begin{tikzpicture}
\begin{axis}[
  width=12cm, height=5.6cm,
  xlabel={number of flips $n$}, ylabel={$\P(A_n)=1-2^{-n}$},
  xmin=0, xmax=10.5, ymin=0, ymax=1.08,
  xtick={1,2,3,4,5,6,7,8,9,10},
  ytick={0,0.25,0.5,0.75,1.0},
  axis lines=left, grid=major, grid style={gray!18},
]
\addplot[dashed, gray, thick] coordinates {(0,1) (10.5,1)};
\addplot[only marks, mark=*, mark size=2pt, distblue]
  coordinates {(1,0.5)(2,0.75)(3,0.875)(4,0.9375)(5,0.96875)
  (6,0.984375)(7,0.9921875)(8,0.99609375)(9,0.998046875)(10,0.999023438)};
\node[anchor=south east, font=\footnotesize, gray] at (axis cs:10.5,1) {$\P(A)=1$};
\end{axis}
\end{tikzpicture}
\caption{Continuity in action for ``eventually heads.'' The finite-stage
probabilities $\P(A_n)=1-2^{-n}$ climb toward the limiting value $\P(A)=1$
(dashed line). Continuity of probability is precisely the guarantee that the
limit of these computable dots equals the probability of the limiting event.}
\label{ch03:fig-continuity}
\end{figure}

\begin{trapbox}
\textbf{Probability zero does not mean impossible (and probability one does not
mean certain).} In Example~\ref{ch03:ex-eventually}, the outcome ``tails forever''
is a perfectly well-defined element of the sample space; it just sits inside a set
of probability zero. In continuous models this is the norm, not the exception:
every \emph{individual} value of a continuous measurement has probability zero,
yet one of them does occur. ``Probability zero'' means \emph{negligible in the
limit / of zero area}, which is a strictly weaker notion than ``cannot happen.''
Conflating the two leads to genuine paradoxes; keeping them separate is a mark of
real understanding.
\end{trapbox}

\section{Event algebras: a gentle preview of $\sigma$-algebras}
\label{ch03:algebra}

We close by returning to the fine print of Definition~\ref{ch03:def-axioms}: the
collection $\mathcal{F}$ of events on which $\P$ is defined. Why not just take
\emph{all} subsets of $\Omega$ and be done with it?

\begin{intuitionbox}
Think about which sets you actually need to be ``events.'' If $A$ is an event,
you certainly want ``not $A$'' to be an event too --- it would be bizarre to be
able to bet on $A$ but not on its negation. If $A$ and $B$ are events, you want
``$A$ or $B$'' and ``$A$ and $B$'' to be events. And from
Section~\ref{ch03:continuity}, if $A_1, A_2, \dots$ are events you want their
\emph{countable} union to be an event, so that limiting statements like
``eventually heads'' are themselves things you can assign a probability to. A
collection of sets closed under exactly these operations is what we need ---
nothing more, nothing less.
\end{intuitionbox}

\begin{definition}[$\sigma$-algebra / event algebra]
\label{ch03:def-sigma-algebra}
A collection $\mathcal{F}$ of subsets of $\Omega$ is a \textbf{$\sigma$-algebra}
(or $\sigma$-field) if:
\begin{enumerate}
\item $\Omega \in \mathcal{F}$;
\item if $A \in \mathcal{F}$ then $A^c \in \mathcal{F}$ (closed under
complements);
\item if $A_1, A_2, \dots \in \mathcal{F}$ then $\bigcup_{n} A_n \in \mathcal{F}$
(closed under countable unions).
\end{enumerate}
The members of $\mathcal{F}$ are exactly the \textbf{events} to which $\P$ assigns
probabilities.
\end{definition}

These three closure properties are deliberately the set-operation mirror of the
three things we want to be able to \emph{say} about events. From them everything
else follows: $\varnothing = \Omega^c \in \mathcal{F}$, and countable
\emph{intersections} are in $\mathcal{F}$ too, because $\bigcap_n A_n =
\big(\bigcup_n A_n^c\big)^c$ by De Morgan's law --- complement, union,
complement, all of which stay inside $\mathcal{F}$.

\begin{example}[Two extreme $\sigma$-algebras, and a useful one]
\label{ch03:ex-sigma}
\begin{itemize}
\item The \emph{trivial} $\sigma$-algebra $\set{\varnothing, \Omega}$: you can
only ask ``did anything happen?'' Useless, but legal.
\item The \emph{power set} of all subsets: the richest possible, and the right
choice whenever $\Omega$ is finite or countable.
\item For one coin flip with $\Omega = \set{H,T}$, the power set is
$\set{\varnothing, \set{H}, \set{T}, \set{H,T}}$ --- four events. Check: it
contains $\Omega$, it is closed under complement ($\set{H}^c = \set{T}$), and
under union. This is a $\sigma$-algebra with $2^2 = 4$ members.
\end{itemize}
\end{example}

\begin{formalbox}
\textbf{Why can't we always use all subsets?} For uncountable $\Omega$, such as
$\Omega = [0,1]$ with ``uniformly random point,'' it is a theorem (relying on the
axiom of choice) that there is \emph{no} way to assign a length to every subset of
$[0,1]$ while keeping countable additivity and translation invariance --- the
so-called non-measurable (Vitali) sets break it. The escape is to define $\P$
only on a $\sigma$-algebra rich enough to contain every set we could ever
construct in practice --- the \emph{Borel} $\sigma$-algebra, generated by the
intervals --- and to simply not ask about the pathological sets. This is not a
cop-out; it is the precise boundary between what can and cannot be measured, and
mapping that boundary is the achievement of measure theory. We develop all of
this carefully in the measure-theoretic chapters; for now, the takeaway is that
$\mathcal{F}$ exists for good reason and that for the finite and countable
problems of the next many chapters you may safely take $\mathcal{F}$ to be all
subsets and forget the subtlety entirely.
\end{formalbox}

\begin{connectionbox}
A $\sigma$-algebra is more than a technical fence --- it encodes
\emph{information}. The events in $\mathcal{F}$ are precisely the questions you are
allowed to answer; a smaller $\sigma$-algebra is a coarser, more ignorant view of
the world. This reading becomes powerful when we study conditional expectation
and martingales, where a growing family of $\sigma$-algebras (a
\emph{filtration}) models information accumulating over time --- what you know at
each moment as a random process unfolds. The humble closure rules of
Definition~\ref{ch03:def-sigma-algebra} are the seed of that entire theory. The
algebraic flavor of ``closed under operations'' also echoes the subspace and
span ideas from the companion linear algebra book: in both cases we ask what a
set of objects generates under the allowed moves.
\end{connectionbox}

\section{Chapter summary}
\label{ch03:summary}

We built the formal stage for all of probability:
\begin{itemize}
\item A random experiment is modeled by a \textbf{sample space} $\Omega$ of
outcomes; \textbf{events} are subsets, and the operations of logic (and, or, not,
implies) become the operations of sets (intersection, union, complement,
inclusion).
\item A \textbf{probability measure} $\P$ obeys three Kolmogorov axioms:
nonnegativity, normalization ($\P(\Omega)=1$), and countable additivity over
disjoint events. The triple $(\Omega, \mathcal{F}, \P)$ is a \textbf{probability
space}.
\item From these three axioms alone we \emph{derived} every elementary rule:
$\P(\varnothing)=0$, the complement rule $\P(A^c)=1-\P(A)$ (so $\P(A)\in[0,1]$),
monotonicity, the difference rule, inclusion--exclusion, and the union bound.
\item When outcomes are \textbf{equally likely}, the axioms force
$\P(A) = \abs{A}/\abs{\Omega}$, turning probability into counting --- but this is
a theorem about a special case, not a definition.
\item \textbf{Continuity of probability} (equivalent to countable additivity)
lets us pass to limits of growing or shrinking events, which is what every
asymptotic theorem in the book depends on, and which forces us to distinguish
\emph{probability zero} from \emph{impossible}.
\item The events form a \textbf{$\sigma$-algebra} $\mathcal{F}$, closed under
complement and countable union --- a structure that is invisible for finite
problems but indispensable for the continuum, and which secretly encodes
\emph{information}.
\end{itemize}

Everything that follows --- conditional probability, random variables,
expectation, the great limit theorems --- is built on top of this small, sturdy
foundation. With the stage set, we can now bring on the actors.

\section{Exercises}
\label{ch03:exercises}

\begin{exercise}[Set translations]
\label{ch03:exr-translate}
Let $A$, $B$, $C$ be events. Express each of the following in set notation, and
then in words describe its complement.
\begin{enumerate}[label=(\alph*)]
\item Exactly one of $A$, $B$ occurs.
\item At least two of $A$, $B$, $C$ occur.
\item None of $A$, $B$, $C$ occurs.
\end{enumerate}
\end{exercise}

\begin{exercise}[Warm-up computations]
\label{ch03:exr-warmup}
Suppose $\P(A) = 0.6$, $\P(B) = 0.5$, and $\P(A \cap B) = 0.3$. Compute
$\P(A \cup B)$, $\P(A^c)$, $\P(A^c \cap B^c)$, $\P(A \cap B^c)$, and
$\P(A^c \cup B^c)$. (Hint: De Morgan turns the last two complements into things
you already have.)
\end{exercise}

\begin{exercise}[A bound becomes an equality]
\label{ch03:exr-bound}
For events $A$ and $B$, prove that $\P(A \cap B) \ge \P(A) + \P(B) - 1$. This is
\emph{Bonferroni's inequality}; show it follows from
$\P(A\cup B)\le 1$ and inclusion--exclusion. When is it an equality?
\end{exercise}

\begin{exercise}[Counting with equally likely outcomes]
\label{ch03:exr-counting}
A fair die is rolled twice. Using the $36$-element equally-likely sample space,
find the probability that (a) the maximum of the two rolls is $5$; (b) the two
rolls differ by exactly $2$; (c) at least one roll is a $6$. For part (c), do it
twice --- directly and via the complement --- and check the answers agree.
\end{exercise}

\begin{exercise}[Why not finite additivity alone?]
\label{ch03:exr-fin-vs-count}
Consider $\Omega = \set{1,2,3,\dots}$ and the geometric assignment
$\P(\set{k}) = 2^{-k}$. (a) Verify $\P(\Omega) = 1$ using countable additivity.
(b) Compute $\P(\set{k : k \text{ even}})$. (c) Explain in a sentence why finite
additivity alone would not even let you conclude $\P(\Omega) = 1$ from the
individual outcome probabilities.
\end{exercise}

\begin{exercise}[Continuity in action]
\label{ch03:exr-continuity}
Let $A_n = $ ``the first head in repeated fair-coin flips occurs on or before
flip $n$.''
\begin{enumerate}[label=(\alph*)]
\item Are the $A_n$ increasing or decreasing? What is their limit event?
\item Compute $\P(A_n)$ and take the limit using
Theorem~\ref{ch03:thm-continuity}.
\item Now let $B_n = $ ``the first $n$ flips are all heads.'' Are these
increasing or decreasing? Find $\bigcap_n B_n$ and its probability, and interpret
the result in light of the trap box on probability zero.
\end{enumerate}
\end{exercise}

\begin{exercise}[Inclusion--exclusion for three]
\label{ch03:exr-ie3}
In a survey, $50\%$ of people own a cat, $40\%$ own a dog, $30\%$ own a fish;
$20\%$ own both a cat and a dog, $15\%$ a cat and a fish, $10\%$ a dog and a fish,
and $5\%$ own all three. Using Theorem~\ref{ch03:thm-ie3}, find the fraction who
own \emph{at least one} of the three pets, and then the fraction who own
\emph{none}.
\end{exercise}

\begin{exercise}[The union bound is sometimes all you can do]
\label{ch03:exr-union-bound}
A server runs $100$ independent jobs; each fails with probability $0.001$. Give
an upper bound on the probability that \emph{at least one} job fails using the
union bound. Then compute the \emph{exact} probability using independence
($1 - 0.999^{100}$) and compare. By roughly how much does the union bound
overshoot, and why is the gap so small here?
\end{exercise}

\begin{exercise}[Proving inclusion--exclusion for three events]
\label{ch03:exr-prove-ie3}
Prove Theorem~\ref{ch03:thm-ie3} from the two-event version. (Hint: write
$A \cup B \cup C = (A \cup B) \cup C$, apply two-event inclusion--exclusion to the
outer union, then distribute $(A\cup B)\cap C = (A\cap C)\cup(B\cap C)$ and apply
it again. Keep careful track of the triple-overlap term.)
\end{exercise}

\begin{exercise}[A $\sigma$-algebra by hand]
\label{ch03:exr-sigma}
Let $\Omega = \set{1,2,3,4}$. (a) Write down the smallest $\sigma$-algebra
containing the single set $\set{1}$. (b) How many members does the power-set
$\sigma$-algebra of $\Omega$ have? (c) Give an example of a collection of subsets
of $\Omega$ that is \emph{not} a $\sigma$-algebra, and state which closure
property fails.
\end{exercise}

\begin{exercise}[Boole, Bonferroni, and a sandwich]
\label{ch03:exr-sandwich}
For events $A_1, \dots, A_n$, the union bound gives an upper estimate of
$\P(\bigcup_i A_i)$. Prove the matching \emph{lower} estimate
$\P(\bigcup_i A_i) \ge \sum_i \P(A_i) - \sum_{i<j}\P(A_i \cap A_j)$. (Hint: this
is the first Bonferroni inequality; induct on $n$ using two-event
inclusion--exclusion and monotonicity.) Together with Boole's inequality this
sandwiches the union probability between two computable bounds.
\end{exercise}

\begin{exercise}[Conceptual: choosing the sample space]
\label{ch03:exr-conceptual}
You flip a fair coin three times. Propose two \emph{different} sample spaces for
this experiment: one in which the outcomes are equally likely, and one in which
they are not. For each, give the sample space explicitly, state $\P$ of one
outcome, and explain which space you would prefer for computing ``the probability
of getting exactly two heads,'' and why.
\end{exercise}
